% ============================================================
%  main_fa.tex — نسخه‌ی فارسی پایان‌نامه
%  قالب: دانشکده مهندسی صنایع، دانشگاه علم و صنعت ایران
%  مستخرج از final-rev07.docx
% ============================================================

\documentclass[12pt,a4paper]{article}

% ============================================================
%  ریاضیات
% ============================================================

\usepackage{amsmath,amssymb,amsthm,mathtools}

% ============================================================
%  رنگ‌ها و جداول
% ============================================================

\usepackage[table]{xcolor}
\usepackage{booktabs}
\usepackage{array}
\usepackage{multirow}
\usepackage{colortbl}
\usepackage{longtable}

% ============================================================
%  تصویر و گرافیک
% ============================================================

\usepackage{graphicx}
\usepackage{float}
\usepackage{caption}
\usepackage{subcaption}
\usepackage{rotating}

% ============================================================
%  TikZ / PGF
% ============================================================

\usepackage{tikz}
\usetikzlibrary{positioning, decorations.pathreplacing, shapes, arrows,
                backgrounds, fit, patterns, shapes.geometric, arrows.meta, calc}
\usepackage{pgfplots}
\pgfplotsset{compat=1.18}

% ============================================================
%  الگوریتم‌ها
% ============================================================

\usepackage{algorithm}
\usepackage{algpseudocode}

% ============================================================
%  چیدمان صفحه  (منطبق با قالب IUST)
%  راست = لبه صحافی (۴ سانت) / چپ = حاشیه بیرونی (۳ سانت)
% ============================================================

\usepackage{geometry}
\geometry{
  a4paper,
  top    = 3cm,
  bottom = 3cm,
  right  = 4cm,   % لبه صحافی — RTL
  left   = 3cm,
  headheight = 14pt,
  headsep    = 0.5cm,
  footskip   = 0.8cm,
}

% ============================================================
%  سرصفحه و پاصفحه
% ============================================================

\usepackage{fancyhdr}
\pagestyle{fancy}
\fancyhf{}
\fancyfoot[C]{{\fontsize{13pt}{16pt}\selectfont\thepage}}
\renewcommand{\headrulewidth}{0.4pt}
\renewcommand{\footrulewidth}{0pt}

% ============================================================
%  فاصله‌گذاری خطوط (منطبق با استایل 'a' در قالب: ۱٫۲ برابر)
% ============================================================

\usepackage{setspace}
\setstretch{1.2}

% ============================================================
%  فاصله‌ی پاراگراف
% ============================================================

\setlength{\parskip}{0pt}
\setlength{\parindent}{1em}

% ============================================================
%  فهرست‌بندی
% ============================================================

\usepackage{enumitem}

% ============================================================
%  سایر بسته‌ها
% ============================================================

\usepackage{mdframed}
\usepackage{listings}
\usepackage{placeins}
\usepackage{appendix}
\usepackage{tabularx}
\usepackage{microtype}

% ============================================================
%  رنگ‌های تگ دسته‌بندی
% ============================================================

\definecolor{TagScheduleBg}{HTML}{EAF3DE}
\definecolor{TagScheduleFg}{HTML}{3B6D11}
\definecolor{TagContractBg}{HTML}{E6F1FB}
\definecolor{TagContractFg}{HTML}{0C447C}
\definecolor{TagCostBg}{HTML}{FAEEDA}
\definecolor{TagCostFg}{HTML}{633806}
\definecolor{TagPenaltyBg}{HTML}{FCEBEB}
\definecolor{TagPenaltyFg}{HTML}{791F1F}
\definecolor{TagTermBg}{HTML}{F0EBFA}
\definecolor{TagTermFg}{HTML}{4B1F8C}
\definecolor{RowAlt}{HTML}{F8F7F4}

% رنگ‌های نمودار PRISMA
\definecolor{prismaBlue}{RGB}{31,119,180}
\definecolor{prismaLightBlue}{RGB}{174,214,241}
\definecolor{prismaGreen}{RGB}{44,160,44}
\definecolor{prismaLightGreen}{RGB}{198,239,206}
\definecolor{prismaGray}{RGB}{127,127,127}
\definecolor{prismaOrange}{RGB}{214,39,40}
\definecolor{prismaLightOrange}{RGB}{255,205,210}

% رنگ‌های جدول مرور ادبیات
\definecolor{symgreen}{HTML}{1A7A4A}
\definecolor{symamber}{HTML}{B45309}
\definecolor{symgrey}{HTML}{9CA3AF}
\definecolor{rarebg}{HTML}{EFF6FF}
\definecolor{headbg}{HTML}{1E3A5F}
\definecolor{headfg}{HTML}{FFFFFF}
\definecolor{pwbg}{HTML}{14532D}
\definecolor{rowodd}{HTML}{F8FAFC}

% ============================================================
%  تِم نمودار جریان
% ============================================================

\colorlet{flowViolet}{violet!70!black}
\colorlet{flowVioletBg}{violet!10}
\colorlet{flowOrange}{orange!70!red}
\colorlet{flowOrangeBg}{orange!8}
\colorlet{flowTeal}{teal!70!black}
\colorlet{flowTealBg}{teal!8}
\colorlet{flowBlue}{blue!60!black}
\colorlet{flowBlueBg}{blue!8}
\colorlet{flowGray}{gray!60}
\colorlet{flowGrayBg}{gray!6}
\colorlet{flowRed}{red!55!black}
\colorlet{flowRedBg}{red!6}

\tikzset{
  flowColorA/.style      = {draw=flowViolet, fill=flowVioletBg},
  flowColorB/.style      = {draw=flowOrange, fill=flowOrangeBg},
  flowColorC/.style      = {draw=flowTeal,   fill=flowTealBg},
  flowColorD/.style      = {draw=flowBlue,   fill=flowBlueBg},
  flowColorGray/.style   = {draw=flowGray,   fill=flowGrayBg},
  flowColorReject/.style = {draw=flowRed,    fill=flowRedBg},
  arrowMain/.style       = {-{Latex[length=6pt]}, thick, black!70},
  arrowReject/.style     = {-{Latex[length=6pt]}, thick, flowRed, dashed},
}

% ============================================================
%  استایل‌های TikZ
% ============================================================

\tikzset{
  box/.style  = {draw, rectangle, rounded corners=4pt,
                 text width=7.0cm, align=center,
                 minimum height=1.2cm, inner sep=7pt,
                 fill=white, draw=black!70},
  excl/.style = {draw, rectangle, rounded corners=4pt,
                 text width=5.5cm, align=left,
                 minimum height=1.2cm, inner sep=7pt,
                 fill=pink!25, draw=red!50!black},
  incl/.style = {draw, rectangle, rounded corners=4pt,
                 text width=7.0cm, align=center,
                 minimum height=1.2cm, inner sep=7pt,
                 fill=green!15, draw=green!60!black},
  phase/.style= {draw, rectangle, rounded corners=4pt,
                 text width=1.5cm, align=center,
                 minimum height=1.2cm, inner sep=5pt,
                 fill=blue!10, draw=blue!50!black,
                 font=\bfseries\small},
  arr/.style  = {->, >=stealth, thick},
  arr2/.style = {->, >=stealth, dashed, thick},
}

% ============================================================
%  ماکروهای تگ دسته‌بندی
% ============================================================

\newcommand{\tagschedule}{%
  \,\colorbox{TagScheduleBg}{\textcolor{TagScheduleFg}{\scriptsize\textsf{زمان‌بندی}}}}
\newcommand{\tagcontract}{%
  \,\colorbox{TagContractBg}{\textcolor{TagContractFg}{\scriptsize\textsf{قرارداد}}}}
\newcommand{\tagcost}{%
  \,\colorbox{TagCostBg}{\textcolor{TagCostFg}{\scriptsize\textsf{هزینه}}}}
\newcommand{\tagpenalty}{%
  \,\colorbox{TagPenaltyBg}{\textcolor{TagPenaltyFg}{\scriptsize\textsf{جریمه}}}}
\newcommand{\tagterm}{%
  \,\colorbox{TagTermBg}{\textcolor{TagTermFg}{\scriptsize\textsf{فسخ}}}}

% ============================================================
%  محیط‌های قضیه
% ============================================================

\newtheorem{theorem}{قضیه}
\newtheorem{lemma}[theorem]{لم}
\newtheorem{proposition}[theorem]{گزاره}
\newtheorem{corollary}[theorem]{نتیجه}
\theoremstyle{definition}
\newtheorem{definition}{تعریف}
\newtheorem{assumption}{فرض}
\theoremstyle{remark}
\newtheorem{remark}{ملاحظه}

% ============================================================
%  دستورات ریاضی سفارشی
% ============================================================

\newcommand{\E}{\mathbb{E}}
\newcommand{\Var}{\text{Var}}
\newcommand{\Cov}{\text{Cov}}
\newcommand{\R}{\mathbb{R}}
\newcommand{\N}{\mathbb{N}}
\newcommand{\argmax}{\operatorname{argmax}}
\newcommand{\argmin}{\operatorname{argmin}}

% ============================================================
%  کد درون‌خطی و بلوک‌های کد
% ============================================================

\usepackage{tcolorbox}
\tcbuselibrary{skins}

\definecolor{inlineBg}{HTML}{F0F0F0}
\newcommand{\code}[1]{%
  \tcbox[on line,
    colback=inlineBg, colframe=inlineBg,
    boxrule=0pt, arc=2pt,
    top=0.5pt, bottom=0.5pt, left=2pt, right=2pt,
    fontupper=\ttfamily\footnotesize\color{black}]{#1}}

\lstset{
    basicstyle=\ttfamily\small\color{black},
    backgroundcolor=\color{inlineBg},
    frame=single,
    framerule=0pt,
    rulecolor=\color{inlineBg},
    breaklines=true,
    keepspaces=true,
    showstringspaces=false,
}

% ============================================================
%  پیوندها و منابع
%  — natbib باید پیش از xepersian بارگذاری شود
% ============================================================

\usepackage{natbib}
\usepackage{hyperref}
\usepackage{cleveref}

\hypersetup{colorlinks=true, linkcolor=blue, citecolor=blue, urlcolor=blue}

% ============================================================
%  فونت‌ها و زبان
%  — xepersian حتماً باید آخرین بسته باشد
%  B Zar  → fc-list: "B Zar"  style=Regular
%  ZarBold → fc-list: "Zar"   style=Bold
% ============================================================

\usepackage{fontspec}
\usepackage{xepersian}

\settextfont[
  Scale    = 1.0,
  BoldFont = Zar,
]{B Zar}

% FIX 1: \setlatintypeface was removed in xepersian v22+.
%         Use \setlatinfont instead.
\setlatinfont[Scale=1.0]{Times New Roman}

% FIX 2: \ZWNJ is only defined after \begin{document} by xepersian.
%         Deferring via \AtBeginDocument ensures \ZWNJ is already available
%         when \zwnj is first used, preventing "Undefined control sequence" errors.
\AtBeginDocument{\providecommand{\zwnj}{\ZWNJ}}

% فونت عنوان‌های بخش
\newfontfamily\headingfont[BoldFont=Zar]{B Zar}

% فونت صفحه عنوان
\newfontfamily\titlefont[BoldFont=Zar]{B Zar}

% ============================================================
%  TikZ را به حالت چپ‌به‌راست (LTR) تنظیم می‌کنیم
%  تا متن انگلیسی داخل نمودارها درست نمایش داده شود
% ============================================================

\AtBeginEnvironment{tikzpicture}{\latintext}
\AtBeginEnvironment{pgfpicture}{\latintext}

% ============================================================
%  عنوان مقاله
% ============================================================

\title{%
  {\titlefont\bfseries\large
   بودجه\zwnj{}بندی سبد پروژه به\zwnj{}عنوان یک فرایند تصمیم\zwnj{}گیری مارکوف:\\[4pt]
   طراحی محیط، ارزیابی خط پایه و تحلیل شکاف بهینگی%
  }
}

\author{%
  \begin{tabular}{c}
    {\headingfont\bfseries مهریار قرقانی}%
      \thanks{نویسنده‌ی مسئول:
        \lr{\href{mailto:mahriargharaghani@gmail.com}{mahriargharaghani@gmail.com}}} \\[2pt]
    {\headingfont\bfseries سید فرید قنادپور} \\[6pt]
    {\itshape\small
      گروه مهندسی صنایع، دانشگاه علم و صنعت ایران} \\
    {\itshape\small تهران، ایران}
  \end{tabular}
}

\date{}

% ============================================================
%  سند
% ============================================================

\begin{document}

\maketitle

% ============================================================
%  چکیده
% ============================================================

\begin{abstract}

تخصیص یک ذخیره‌ی سرمایه‌ی مشترک در میان سبدی از پروژه‌های قراردادی در حال اجرا،
یک تصمیم\zwnj{}گیری پرتکرار و حساس برای پیمانکاران مهندسی است.
رویکردهای کلاسیک یا آینده‌ی کامل را مفروض می\zwnj{}گیرند یا سبد را به یک پروژه‌ی واحد
تقلیل می\zwnj{}دهند و در نتیجه هیچ ابزار اصولی‌ای برای مسئله‌ی ترتیبی و چندپروژه‌ای
در شرایط عدم\zwnj{}قطعیت دنیای واقعی در اختیار پیشه\zwnj{}ورزان نمی\zwnj{}گذارند.
این مقاله آن شکاف را با ارائه‌ی یک محیط شبیه\zwnj{}سازی کامل و تأییدشده
پر می\zwnj{}کند که این مسئله را برای روش\zwnj{}های داده\zwnj{}محور قابل حل می\zwnj{}سازد.

مسئله را به\zwnj{}صورت یک فرایند تصمیم\zwnj{}گیری مارکوف مقیدشده (\lr{CMDP}) صورت\zwnj{}بندی
می\zwnj{}کنیم که در آن مدیر سبد در هر دوره‌ی گسسته کسرهای بودجه را میان
$N$ پروژه‌ی موازی تخصیص می\zwnj{}دهد.
حالت، شاخص\zwnj{}های مدیریت ارزش کسب\zwnj{}شده (\lr{EVM})،
پیشرفت برنامه\zwnj{}ریزی\zwnj{}شده‌ی منحنی\lr{-S}،
پیش\zwnj{}نگاه پرداخت\zwnj{}های نقطه‌ی عطف، و یک شمارنده‌ی تحمل نقض
را برای هر پروژه رمزگذاری می\zwnj{}کند.
دینامیک گذار چرخه‌ی کامل جریان نقدی قرارداد
\lr{FIDIC} را بازتولید می\zwnj{}کند:
پرداخت پیش\zwnj{}پرداخت و بازیابی آن، پرداخت\zwnj{}های ناخالص نقطه‌ی عطف
متناسب با پیشرفت، کسر و آزادسازی حسن انجام کار،
هزینه‌ی بهره بر برداشت از خزانه، و دروازه‌ی دوشرطی فسخ قرارداد.
سیگنال پاداش، جریان نقدی خالص سبد با تنزیل $\gamma$ است که
به \lr{NPV} یک اپیزود انباشته می\zwnj{}شود و معیار اصلی عملکرد را تشکیل می\zwnj{}دهد.

محیط به\zwnj{}صورت یک بسته‌ی پیکربندی\zwnj{}پذیر سازگار با \lr{Gymnasium}
منتشر می\zwnj{}شود که دارای یک دفتر کل \lr{SQLite} است
و هر تخصیص، رویداد جریان نقدی و پرچم نقض را ثبت می\zwnj{}کند---
این ویژگی مشکل جعبه\zwnj{}سیاه ذاتی سیاست\zwnj{}های یادگیری\zwnj{}شده را حذف می\zwnj{}کند.
برای ایجاد یک خط پایه‌ی عملکرد و انگیزه\zwnj{}بخشی به
استفاده از عامل\zwnj{}های یادگیری\zwnj{}شده،
سه اکتشافی قطعی حریصانه (\lr{Reach-Plan}، \lr{Catchup-T} و \lr{Catchup-Next-T})
و یک بهینه‌ی اطلاعات\zwnj{}کامل محاسبه\zwnj{}شده‌ی دستی را
روی چهار پیکربندی معیار با پیچیدگی سبد فزاینده اجرا می\zwnj{}کنیم.
نتایج یک شکاف سیستماتیک و روبه\zwnj{}افزایش میان خطوط پایه‌ی حریصانه
و راه\zwnj{}حل بهینه با افزایش پیچیدگی سبد را نشان می\zwnj{}دهند
و تأیید می\zwnj{}کنند که تخصیص مبتنی بر قوانین استاندارد ناکافی است
و یک هدف بهبود مشخص برای عامل\zwnj{}های یادگیری تقویتی وجود دارد.
آموزش چنین عامل\zwnj{}هایی و ارزیابی آن\zwnj{}ها در برابر یک کران
برنامه‌ریزی خطی رسمی به\zwnj{}عنوان گام\zwnj{}های طبیعی بعدی
شناسایی می\zwnj{}شوند و در کارهای بعدی به آن\zwnj{}ها پرداخته خواهد شد.

\medskip
\noindent\textbf{کلیدواژه\zwnj{}ها:}
مدیریت سبد پروژه؛
تخصیص بودجه تحت عدم\zwnj{}قطعیت عملکرد؛
فرایند تصمیم\zwnj{}گیری مارکوف مقیدشده؛
یادگیری تقویتی؛
قراردادهای \lr{FIDIC}؛
مدیریت ارزش کسب\zwnj{}شده؛
محیط \lr{Gymnasium}؛
خط پایه‌ی حریصانه؛
شکاف بهینگی.

\end{abstract}

\FloatBarrier

% ============================================================
%  مقدمه
% ============================================================

\section{مقدمه}
\label{sec:introduction}
\input{sections/01_introduction_fa.tex}

\FloatBarrier

% ============================================================
%  مرور ادبیات
% ============================================================

\section{مرور ادبیات}
\label{sec:literature}
\input{sections/02_literature_fa.tex}

\FloatBarrier

% ============================================================
%  صورت‌بندی مسئله
% ============================================================

% \section{صورت‌بندی مسئله}
% \label{sec:problem}
% \input{sections/03_problem_formulation_fa.tex}

% \FloatBarrier

% % ============================================================
% %  ارزیابی
% % ============================================================

% \section{ارزیابی}
% \label{sec:evaluation}
% \input{sections/04_evaluation_fa.tex}

% \FloatBarrier

% % ============================================================
% %  نتیجه‌گیری
% % ============================================================

% \section{نتیجه‌گیری}
% \label{sec:conclusion}
% \input{sections/05_conclusion_fa.tex}

% \FloatBarrier

% ============================================================
%  منابع
%  — مسیر کامل references.bib که کنار main_fa.tex است
% ============================================================

\bibliographystyle{apalike}
\bibliography{references}
\clearpage

% ============================================================
%  پیوست الف — نمودار جریان PRISMA 2020
% ============================================================

% \appendix
% \section{پیوست الف: نمودار جریان \lr{PRISMA 2020}}
% \label{app:prisma}
% \input{sections/appendix_A_fa.tex}
% \FloatBarrier

% % ============================================================
% %  پیوست ب — گواهی سازگاری Gymnasium
% % ============================================================

% \section{پیوست ب: گواهی سازگاری \lr{Gymnasium}}
% \label{app:gym_certificate}
% \input{sections/appendix_B_fa.tex}
% \FloatBarrier

% % ============================================================
% %  پیوست ج — معماری مخزن و نقشه‌ی مفهومی
% % ============================================================

% \section{پیوست ج: معماری مخزن و نقشه‌ی مفهومی}
% \label{app:repo_map}
% \input{sections/appendix_C_fa.tex}

\end{document}